Two benchmark functions are treated as expensive simulations to test how far a cheap surrogate can be trusted. For the Rosenbrock function (common) and the two-variable Trid function (assigned), 60 evaluations from a maximin Latin hypercube design are split into 45 training and 15 held-out test points. Three surrogates are compared: polynomial regression with its degree chosen by five-fold cross-validation on the training set only, a cubic radial basis function (RBF) interpolant, and a Gaussian process (GP). They are assessed with test metrics, side-by-side error maps, learning curves, surrogate-based optimization with verification, and a reduced-box extrapolation test. Cross-validation recovers the true polynomial degrees (4 and 2), and the fitted coefficients match the exact expansions to within $4\times10^{-13}$. The GP reaches a test normalized root-mean-square error of $7.8\times10^{-6}$ and $1.0\times10^{-6}$ and places the two minima within 0.013 and $6\times10^{-4}$ of the truth. The RBF fails at the domain boundaries and creates a spurious negative minimum in the Rosenbrock valley, which only verification exposes. Outside a reduced training box, errors grow by factors of up to 2300. The hand-derived Trid minimum, $f=-2$ at $(2,2)$, is confirmed analytically and numerically.
